\SourceLeaf{321}%
oder einer anderen Variante (unter den vier jeweils zusammengehörigen), erkennt man erst hinterher,\\\hspace*{2em}wenn man nicht gleich zu Anfang mein Kriterium angewandt hat.

Ich gebe das Kriterium selbst hier nicht an, um Ihnen nicht das Vergnügen zu nehmen, es selbst zu finden. Jedenfalls hoffe ich, durch das Vorstehende hinreichend deutlich gemacht zu haben, \textls[120]{worum es sich handelt}. Übrigens ist es wohl unnötig, zu bemerken, dass der ganze Gegenstand höchst elementar ist. Umso erstaunlicher ist es, dass man ihn so völlig übersehen hat.

\mbox{. . . . . . . . . .}

\Needspace{12\baselineskip}\begin{center}[2.]\end{center}

\begin{center}\textsc{Gauss} an \textsc{Gerling}. Göttingen, 24. Oktober 1840.\end{center}

\indent \mbox{. . . . . . . .} Ihre kleine Schrift über die \textsc{Pothenot}sche Aufgabe habe ich mit Vergnügen gelesen. Sie wird für alle, die die Aufgabe praktisch anwenden wollen, sehr lehrreich und nützlich sein. Ich habe mich aber oft darüber gewundert, dass die vielen Autoren, die sich damit beschäftigt haben, einen in theoretischer Hinsicht ganz wesentlichen Teil völlig übergangen haben, nämlich die Angabe der mathematischen Bedingungen, unter denen die Aufgabe \textls[120]{physisch möglich} ist, und die Darstellung dieser Bedingungen in einer einfachen, eleganten Form. Ich mache mich wohl folgendermaßen am leichtesten verständlich.

Wendet man die Rechenvorschriften auf vorgegebene numerische Daten an, so erhält man immer ein einziges bestimmtes Ergebnis (es sei denn, die vier Punkte liegen auf einem Kreis, dann erhält man $\frac{0}{0}$). Man erkennt aber leicht, dass man \textls[120]{dasselbe} Ergebnis erhält, ob nun verlangt wird, dass der zweite Punkt um den Winkel $A$ rechts vom ersten liegt oder um $180^0-A$ links von ihm. Ebenso verhält es sich beim dritten Punkt. Anders ausgedrückt gilt Folgendes.

Ist vorgeschrieben, dass der zweite Punkt um $A$ rechts vom ersten liegt und der dritte um $B$ rechts vom ersten, so findet man \textls[120]{dasselbe} Ergebnis, das man bei drei anderen Datensätzen finden würde, nämlich
